% !TEX program = xelatex
\documentclass[11pt,a4paper]{ctexart}

\usepackage{amsmath,amssymb,amsfonts}
\usepackage{booktabs}
\usepackage{geometry}
\usepackage{hyperref}
\usepackage{enumitem}
\usepackage{xcolor}
\usepackage{listings}
\usepackage{longtable}

\geometry{left=2.5cm,right=2.5cm,top=2.5cm,bottom=2.5cm}
\hypersetup{colorlinks=true,linkcolor=blue,urlcolor=blue}

\lstset{
  basicstyle=\ttfamily\small,
  breaklines=true,
  frame=single,
  backgroundcolor=\color{gray!8},
  columns=fullflexible,
  keepspaces=true,
  showstringspaces=false
}

\title{%
  \textbf{星闪 SLB 物理层}\\[0.4em]
  \Large 6.5.6 调制方式技术文档\\[0.3em]
  \large QPSK \textbullet\ 16QAM \textbullet\ 64QAM \textbullet\ 256QAM \textbullet\ 1024QAM \textbullet\ 4096QAM
}
\author{依据 T/XS 10001-2025《星闪无线通信系统 接入层\\
同步低时延宽带空口 SLB 技术要求和测试方法》整理\\
（团体标准 20250326 版）\\[0.5em]
代码文件：\texttt{slb\_phy\_mod.h} / \texttt{slb\_phy\_mod.c}}
\date{}

\begin{document}
\maketitle
\tableofcontents
\newpage

%======================================================================
\section{协议章节对应}
%======================================================================

本节对应星闪 SLB 接入层物理层协议 T/XS 10001-2025 第 6.5 节``共性基础算法''中调制方式子节，具体涵盖：
\begin{itemize}[nosep]
  \item \textbf{6.5.6.1} QPSK——每 2 比特映射 1 个复数符号，归一化因子 $1/\sqrt{2}$；
  \item \textbf{6.5.6.2} 16QAM——每 4 比特映射 1 符号，归一化因子 $1/\sqrt{10}$；
  \item \textbf{6.5.6.3} 64QAM——每 6 比特映射 1 符号，归一化因子 $1/\sqrt{42}$；
  \item \textbf{6.5.6.4} 256QAM——每 8 比特映射 1 符号，归一化因子 $1/\sqrt{170}$；
  \item \textbf{6.5.6.5} 1024QAM——每 10 比特映射 1 符号，归一化因子 $1/\sqrt{682}$；
  \item \textbf{6.5.6.6} 4096QAM——每 12 比特映射 1 符号，归一化因子 $1/\sqrt{2730}$。
\end{itemize}

%======================================================================
\section{功能定义}
%======================================================================

本模块为 SLB 120\,kHz 物理层中控制信息与数据 TB 提供\textbf{比特域 $\rightarrow$ 符号域}的统一调制算法库。6.5.6 不规定信道编码、加扰或资源映射，只规定\textbf{怎么把}编码/加扰后的比特序列 $b(k)$ 映射为复数调制符号 $d(i)$。

输入通常为 6.5.7 加扰输出 $\tilde{g}_k$ 或极化/LDPC 编码后比特流；输出 $d(i)$ 供 6.2.3 资源映射填入 OFDM RE。调制阶数由 MCS 表（6.2.5）或控制信息固定配置（多数为 QPSK）决定。

\begin{table}[h]
\centering
\small
\begin{tabular}{lllll}
\toprule
\textbf{节号} & \textbf{名称} & \textbf{输入} & \textbf{输出} & \textbf{典型链路用途} \\
\midrule
6.5.6.1 & QPSK & $b(2i),b(2i{+}1)$ & $d(i)\in\mathbb{C}$ & 同步/广播/接入/GCI/TCI \\
6.5.6.2 & 16QAM & $b(4i)\ldots b(4i{+}3)$ & $d(i)$ & 数据 MCS、二阶段 HARQ-ACK \\
6.5.6.3 & 64QAM & $b(6i)\ldots b(6i{+}5)$ & $d(i)$ & 第一类/第二类数据 \\
6.5.6.4 & 256QAM & $b(8i)\ldots b(8i{+}7)$ & $d(i)$ & 高阶数据 MCS \\
6.5.6.5 & 1024QAM & $b(10i)\ldots b(10i{+}9)$ & $d(i)$ & 高阶数据 MCS \\
6.5.6.6 & 4096QAM & $b(12i)\ldots b(12i{+}11)$ & $d(i)$ & 极良好信道数据 \\
\bottomrule
\end{tabular}
\end{table}

\textbf{与 6.5.3 的边界}：6.5.3 伪随机 QPSK 用于 SRS、DMRS 等\textbf{参考信号}（Gold 序列驱动，系数为 $+1/\sqrt{2}$）；6.5.6 用于\textbf{业务比特}调制（含 leading 负号）。二者数学形式相近，接口与比特来源不得混用。

%======================================================================
\section{模块边界（上下游及职责划分）}
%======================================================================

\subsection{上游模块（输入来源）}

\begin{table}[h]
\centering
\small
\begin{tabular}{p{4.2cm}p{4.5cm}p{5.5cm}}
\toprule
\textbf{上游模块} & \textbf{输入给本模块的内容} & \textbf{说明} \\
\midrule
CRC/信道编码（6.5.1/6.2.6） & 编码比特 $g_k$ & 控制信息极化编码后比特 \\
比特加扰（6.5.7） & 加扰比特 $\tilde{g}_k$ & 数据 TB 调制前须完成加扰 \\
MCS/调度配置 & 调制方式枚举 & 表 8\textasciitilde 10 或固定 QPSK \\
PHY 比特缓冲 & $b(0)\ldots b(N_{\mathrm{bit}}{-}1)$ & 按自然顺序连续消费 \\
\bottomrule
\end{tabular}
\end{table}

\subsection{下游模块（输出去向）}

\begin{table}[h]
\centering
\small
\begin{tabular}{p{4.2cm}p{4.5cm}p{5.5cm}}
\toprule
\textbf{下游模块} & \textbf{本模块输出} & \textbf{说明} \\
\midrule
资源映射（6.2.3） & 复数符号 $d(i)$ 序列 & 映射到子载波 RE \\
功率控制（6.5.4） & 符号幅度（归一化） & 平均符号能量为 1，再叠功控 \\
接收解调 & 星座点参考 & 软判决 LLR 须用相同 Gray 映射 \\
\bottomrule
\end{tabular}
\end{table}

\subsection{本模块不负责的内容}

\begin{table}[h]
\centering
\small
\begin{tabular}{p{5.5cm}p{8.5cm}}
\toprule
\textbf{不负责内容} & \textbf{原因} \\
\midrule
信道编码与 CRC & 由 6.5.1、6.2.6 负责 \\
比特加扰 & 由 6.5.7 负责，数据调制输入为 $\tilde{g}_k$ \\
参考信号 QPSK（6.5.3） & Gold 序列驱动，非业务比特 \\
资源网格位置与时频分配 & 由 6.2.3 负责 \\
MCS 选择与链路自适应 & 由 MAC/调度根据 SINR 决定 \\
\bottomrule
\end{tabular}
\end{table}

%======================================================================
\section{在 SLB 通信流程中的角色}
%======================================================================

\begin{itemize}
  \item \textbf{控制信息（6.2.4）}：极化编码 $\rightarrow$ \textbf{6.5.6.1 QPSK} $\rightarrow$ 资源映射；
  \item \textbf{数据 TB（6.2.5）}：CRC $\rightarrow$ LDPC/极化 $\rightarrow$ 6.5.7 加扰 $\rightarrow$ \textbf{6.5.6.x} $\rightarrow$ 资源映射 $\rightarrow$ 6.5.4 功控；
  \item \textbf{TCI 二阶段 HARQ-ACK}：依 G 链路数据调制，选 QPSK 或 16QAM（6.5.6.1/6.5.6.2）。
\end{itemize}

\textbf{推荐复用策略}：工程上维护一份 \texttt{slb\_phy\_mod} 库；按 \texttt{slb\_mod\_type\_e} 分发至各子节实现；QPSK/16QAM 可用查表，64QAM 及以上可用公式或预生成星座表。

%======================================================================
\section{强制要求与关键参数}
%======================================================================

\subsection{通用符号约定}

\begin{table}[h]
\centering
\small
\begin{tabular}{llp{7cm}}
\toprule
\textbf{符号} & \textbf{取值/单位} & \textbf{含义} \\
\midrule
$b(k)$ & $\{0,1\}$ & 输入比特，按自然顺序从编码/加扰流依次取用 \\
$d(i)$ & $\mathbb{C}$ & 第 $i$ 个调制符号（同相 I + 正交 Q） \\
$m$ & 2/4/6/8/10/12 & 每符号比特数 \\
归一化 & $1/\sqrt{N_{\mathrm{avg}}}$ & 保证平均符号能量为 1 \\
leading ``$-$'' & 必选 & 所有 6.5.6.x 公式均含整体负号 \\
\bottomrule
\end{tabular}
\end{table}

\subsection{6.5.6.1 QPSK}

$b(2i)$、$b(2i{+}1)$ 映射为：
\begin{equation}
  d(i) = -\frac{1}{\sqrt{2}}\Bigl\{
    \bigl(2b(2i)-1\bigr) + j\bigl(2b(2i+1)-1\bigr)
  \Bigr\}
  \label{eq:qpsk}
\end{equation}

\subsection{6.5.6.2 16QAM}

$b(4i)\ldots b(4i{+}3)$ 映射为：
\begin{equation}
  d(i) = -\frac{1}{\sqrt{10}}\Bigl\{
    \bigl(2b(4i)-1\bigr)\bigl[2-(1-2b(4i+2))\bigr]
    + j\bigl(2b(4i+1)-1\bigr)\bigl[2-(1-2b(4i+3))\bigr]
  \Bigr\}
  \label{eq:16qam}
\end{equation}

I/Q 各 2 比特，电平 $\{\pm1,\pm3\}$；内层 $2-(1-2b)$ 为 Gray 幅度层级。

\subsection{6.5.6.3 64QAM}

$b(6i)\ldots b(6i{+}5)$，归一化 $1/\sqrt{42}$；I/Q 各 3 比特，电平 $\{\pm1,\pm3,\pm5,\pm7\}$，嵌套结构与 16QAM 相同（三层 $2-(1-2b)$ 组合）。

\subsection{6.5.6.4\textasciitilde 6.5.6.6 高阶 QAM}

\begin{table}[h]
\centering
\small
\begin{tabular}{llll}
\toprule
\textbf{子节} & \textbf{每符号比特} & \textbf{I/Q 比特数} & \textbf{归一化因子} \\
\midrule
256QAM & 8 & 4 & $1/\sqrt{170}$ \\
1024QAM & 10 & 5 & $1/\sqrt{682}$ \\
4096QAM & 12 & 6 & $1/\sqrt{2730}$ \\
\bottomrule
\end{tabular}
\end{table}

嵌套规律：最内层 $A_0=2-(1-2b)\in\{1,3\}$；向外 $A_k=2^{k+1}-(1-2b_k)\cdot A_{k-1}$；实部 $=(2b_{\mathrm{sign,I}}-1)\cdot A_{m/2-1}$，虚部对称；整体乘 leading ``$-$'' 与归一化因子。

\subsection{标准算例（QPSK）}

$b(0)=1,b(1)=0$：括号内 $1-j$，$d(0)=-\frac{1}{\sqrt{2}}(1-j)=\frac{-1+j}{\sqrt{2}}$。

\subsection{6.5.3 与 6.5.6.1 差异}

\begin{table}[h]
\centering
\small
\begin{tabular}{lll}
\toprule
 & \textbf{6.5.3 伪随机 QPSK} & \textbf{6.5.6.1 QPSK} \\
\midrule
比特来源 & Gold $c(2m),c(2m{+}1)$ & 编码/加扰 $b(2i),b(2i{+}1)$ \\
系数 & $+1/\sqrt{2}$ & $-1/\sqrt{2}$（式\eqref{eq:qpsk}） \\
用途 & SRS、DMRS、CSI-RS 等 & 控制信息、数据 TB \\
子载波 & 固定 161（$k{=}80$ 置零） & 由调度决定 \\
\bottomrule
\end{tabular}
\end{table}

%======================================================================
\section{数据类型、长度/维度}
%======================================================================

\subsection{调制类型枚举 \texttt{slb\_mod\_type\_e}}

\begin{table}[h]
\centering
\small
\begin{tabular}{lll}
\toprule
\textbf{枚举值} & \textbf{标准子节} & \textbf{每符号比特 $m$} \\
\midrule
\texttt{SLB\_MOD\_QPSK} & 6.5.6.1 & 2 \\
\texttt{SLB\_MOD\_16QAM} & 6.5.6.2 & 4 \\
\texttt{SLB\_MOD\_64QAM} & 6.5.6.3 & 6 \\
\texttt{SLB\_MOD\_256QAM} & 6.5.6.4 & 8 \\
\texttt{SLB\_MOD\_1024QAM} & 6.5.6.5 & 10 \\
\texttt{SLB\_MOD\_4096QAM} & 6.5.6.6 & 12 \\
\bottomrule
\end{tabular}
\end{table}

\subsection{通用调制接口 \texttt{slb\_mod\_map}}

\begin{table}[h]
\centering
\small
\begin{tabular}{lllll}
\toprule
\textbf{变量} & \textbf{方向} & \textbf{类型} & \textbf{长度} & \textbf{说明} \\
\midrule
\texttt{bits[]} & 输入 & \texttt{const uint8\_t *} & $N_{\mathrm{bit}}$ & 比特缓冲，取 $\{0,1\}$ \\
\texttt{N\_bit} & 输入 & \texttt{uint32\_t} & 1 & 输入比特总数 \\
\texttt{mod\_type} & 输入 & \texttt{slb\_mod\_type\_e} & 1 & 调制方式 \\
\texttt{symbols[]} & 输出 & \texttt{slb\_cplx\_t *} & $N_{\mathrm{bit}}/m$ & 复数符号（float I/Q） \\
\texttt{N\_sym} & 输出 & \texttt{uint32\_t *} & 1 & 实际输出符号数 \\
\bottomrule
\end{tabular}
\end{table}

\texttt{slb\_cplx\_t}：\texttt{struct \{ float re; float im; \}}。

\subsection{单符号接口（各子节）}

\begin{table}[h]
\centering
\small
\begin{tabular}{llll}
\toprule
\textbf{接口} & \textbf{输入} & \textbf{输出} & \textbf{说明} \\
\midrule
\texttt{slb\_mod\_qpsk\_one} & $b(2i),b(2i{+}1)$ & $d(i)$ & 6.5.6.1 \\
\texttt{slb\_mod\_16qam\_one} & $b(4i)\ldots b(4i{+}3)$ & $d(i)$ & 6.5.6.2 \\
\texttt{slb\_mod\_64qam\_one} & 6 比特 & $d(i)$ & 6.5.6.3 \\
\texttt{slb\_mod\_256qam\_one} & 8 比特 & $d(i)$ & 6.5.6.4 \\
\texttt{slb\_mod\_1024qam\_one} & 10 比特 & $d(i)$ & 6.5.6.5 \\
\texttt{slb\_mod\_4096qam\_one} & 12 比特 & $d(i)$ & 6.5.6.6 \\
\bottomrule
\end{tabular}
\end{table}

%======================================================================
\section{时序关系}
%======================================================================

\begin{itemize}[nosep]
  \item \textbf{控制信息}：极化编码完成、比特序确定后、资源映射前调用 QPSK 调制；
  \item \textbf{数据 TB}：6.5.7 加扰完成后、资源映射前调用；调制阶数与 MCS 同步切换；
  \item \textbf{比特消费}：$b(2i)$ 必须先于 $b(2i{+}1)$；跨符号不得乱序；
  \item \textbf{与功控}：调制输出归一化符号后，由 6.5.4 按 $M$、$P_o$、$P_L$ 决定 RE 发射功率。
\end{itemize}

%======================================================================
\section{接口表格（明确的输入/输出）}
%======================================================================

\subsection{\texttt{slb\_mod\_map}}

\begin{longtable}{p{2.5cm}p{2.5cm}p{1.5cm}p{1.5cm}p{2.5cm}p{4cm}}
\toprule
\textbf{参数} & \textbf{类型} & \textbf{长度} & \textbf{必需} & \textbf{来源} & \textbf{说明} \\
\midrule
\endhead
\texttt{bits} & \texttt{const uint8\_t *} & $N_{\mathrm{bit}}$ & 是 & 6.5.7/6.2.6 & 输入比特流 \\
\texttt{N\_bit} & \texttt{uint32\_t} & 1 & 是 & 调度 & 须为 $m$ 的整数倍 \\
\texttt{mod\_type} & \texttt{slb\_mod\_type\_e} & 1 & 是 & MCS/6.2 & 调制方式 \\
\texttt{symbols} & \texttt{slb\_cplx\_t *} & $\geq N_{\mathrm{bit}}/m$ & 是 & --- & 输出缓冲 \\
\texttt{N\_sym} & \texttt{uint32\_t *} & 1 & 是 & --- & 输出符号数 \\
\bottomrule
\end{longtable}

\subsection{\texttt{slb\_mod\_qpsk\_one}}

输入：\texttt{b0}、\texttt{b1} $\in\{0,1\}$；输出：\texttt{slb\_cplx\_t}（式\eqref{eq:qpsk}）。

%======================================================================
\section{处理步骤}
%======================================================================

\subsection{批量调制（通用）}

\begin{enumerate}
  \item 根据 \texttt{mod\_type} 查每符号比特数 $m$；
  \item 校验 $N_{\mathrm{bit}} \bmod m = 0$，否则返回 \texttt{SLB\_PHY\_ERR\_PARAM}；
  \item 对 $i=0\ldots N_{\mathrm{bit}}/m-1$，取 $b(mi)\ldots b(mi{+}m{-}1)$；
  \item 调用对应 \texttt{slb\_mod\_*\_one} 得 $d(i)$；
  \item 写入 \texttt{symbols[i]}，\texttt{*N\_sym} $= N_{\mathrm{bit}}/m$。
\end{enumerate}

\subsection{6.5.6.2 16QAM 单符号（公式拆解）}

\begin{enumerate}
  \item 内层：$A_I=2-(1-2b(4i{+}2))$，$A_Q=2-(1-2b(4i{+}3))$ $\in\{1,3\}$；
  \item 外层符号：$\mathrm{Re}=(2b(4i)-1)\cdot A_I$，$\mathrm{Im}=(2b(4i{+}1)-1)\cdot A_Q$；
  \item $d(i)=-\frac{1}{\sqrt{10}}(\mathrm{Re}+j\mathrm{Im})$。
\end{enumerate}

\subsection{64QAM 及以上}

从最内层 $2-(1-2b)$ 向外递推幅度层级，再乘符号位，最后整体乘 $-1/\sqrt{N_{\mathrm{avg}}}$。高阶建议查表或代码生成，避免手写嵌套错误。

%======================================================================
\section{状态机与时序}
%======================================================================

调制模块为\textbf{无状态}纯函数：每次调用仅依赖当前比特切片与 \texttt{mod\_type}。在 PHY 发送流水线中的位置：

\begin{itemize}[nosep]
  \item 态 1（编码完成）：比特缓冲就绪；
  \item 态 2（加扰完成，数据链路）：$\tilde{g}_k$ 作为 $b(k)$；
  \item 态 3（调制）：$b(k)\to d(i)$；
  \item 态 4（映射/功控）：$d(i)$ 填入 RE 并叠发射功率。
\end{itemize}

%======================================================================
\section{协议中允许本模块改变的参数}
%======================================================================

\begin{table}[h]
\centering
\small
\begin{tabular}{lll}
\toprule
\textbf{参数} & \textbf{来源} & \textbf{影响} \\
\midrule
\texttt{mod\_type} & MCS 表/6.2 固定配置 & 映射公式与频谱效率 \\
输入比特序 & 6.5.7/6.2.6 & 须与收发一致 \\
浮点/定点 Q 格式 & 工程实现 & EVM 与 bit-exact 测试 \\
\bottomrule
\end{tabular}
\end{table}

PHY \textbf{不得}自行升级调制阶数；仅执行上层指定的 6.5.6.x 子节。

%======================================================================
\section{约束与极限条件}
%======================================================================

\begin{table}[h]
\centering
\small
\begin{tabular}{p{4.5cm}p{9.5cm}}
\toprule
\textbf{约束} & \textbf{处理建议} \\
\midrule
$N_{\mathrm{bit}}$ 须为 $m$ 整数倍 & 否则返回 \texttt{SLB\_PHY\_ERR\_PARAM} \\
比特须先加扰（数据链路） & 输入应为 $\tilde{g}_k$，非原始 $g_k$ \\
leading 负号必须实现 & 省略或与 6.5.3 混用将导致 EVM 失败 \\
不得与 6.5.3 接口混用 & 参考信号走 \texttt{slb\_qpsk\_rs\_gen} \\
输出缓冲 $\geq N_{\mathrm{bit}}/m$ & 不足返回 \texttt{SLB\_PHY\_ERR\_BUF} \\
解调 LLR 须同 Gray 映射 & 嵌套 $2-(1-2b)$ 即 Gray 层次 \\
\bottomrule
\end{tabular}
\end{table}

%======================================================================
\section{链路调用对照表}
%======================================================================

\begin{table}[h]
\centering
\small
\begin{tabular}{lll}
\toprule
\textbf{链路/信息块} & \textbf{标准位置} & \textbf{6.5.6 调制} \\
\midrule
同步信息 & 6.2.4.1.2 & QPSK \\
广播信息 & 6.2.4.3.2 & QPSK \\
接入信息 & 6.2.4.4.2 & QPSK \\
GCI & 6.2.4.6.8 & QPSK \\
TCI（含二阶段 HARQ-ACK） & 6.2.4.10.5 & QPSK / 16QAM \\
第一类数据 & 6.2.5.2.2 & QPSK\textasciitilde 4096QAM（MCS 表 8） \\
第二类数据 & 6.2.5.3.2 & QPSK\textasciitilde 4096QAM（表 9/10） \\
\bottomrule
\end{tabular}
\end{table}

%======================================================================
\section{调制方式速查}
%======================================================================

\begin{table}[h]
\centering
\small
\begin{tabular}{llll}
\toprule
\textbf{调制} & \textbf{$m$} & \textbf{归一化} & \textbf{典型场景} \\
\midrule
QPSK & 2 & $1/\sqrt{2}$ & 控制信道、低 MCS 数据 \\
16QAM & 4 & $1/\sqrt{10}$ & 中 MCS、二阶段 ACK \\
64QAM & 6 & $1/\sqrt{42}$ & 中高 MCS 数据 \\
256QAM & 8 & $1/\sqrt{170}$ & 高 MCS 数据 \\
1024QAM & 10 & $1/\sqrt{682}$ & 极高 MCS \\
4096QAM & 12 & $1/\sqrt{2730}$ & 极好信道数据 \\
\bottomrule
\end{tabular}
\end{table}

%======================================================================
\section{API 速查}
%======================================================================

\begin{lstlisting}[language=C]
typedef enum {
    SLB_MOD_QPSK = 0,
    SLB_MOD_16QAM,
    SLB_MOD_64QAM,
    SLB_MOD_256QAM,
    SLB_MOD_1024QAM,
    SLB_MOD_4096QAM
} slb_mod_type_e;

typedef struct { float re; float im; } slb_cplx_t;

/* 批量调制 */
int slb_mod_map(const uint8_t *bits, uint32_t N_bit,
                slb_mod_type_e mod_type,
                slb_cplx_t *symbols, uint32_t *N_sym);

/* 单符号（各子节） */
slb_cplx_t slb_mod_qpsk_one(uint8_t b0, uint8_t b1);
slb_cplx_t slb_mod_16qam_one(const uint8_t b[4]);
slb_cplx_t slb_mod_64qam_one(const uint8_t b[6]);
slb_cplx_t slb_mod_256qam_one(const uint8_t b[8]);
slb_cplx_t slb_mod_1024qam_one(const uint8_t b[10]);
slb_cplx_t slb_mod_4096qam_one(const uint8_t b[12]);

uint8_t slb_mod_bits_per_symbol(slb_mod_type_e mod_type);
\end{lstlisting}

%======================================================================
\section{错误码}
%======================================================================

\begin{table}[h]
\centering
\begin{tabular}{cll}
\toprule
\textbf{宏} & \textbf{值} & \textbf{含义} \\
\midrule
\texttt{SLB\_PHY\_OK} & 0 & 成功 \\
\texttt{SLB\_PHY\_ERR\_PARAM} & $-1$ & 空指针、$N_{\mathrm{bit}}$ 非 $m$ 倍数、非法 \texttt{mod\_type} \\
\texttt{SLB\_PHY\_ERR\_BUF} & $-2$ & 输出符号缓冲区不足 \\
\bottomrule
\end{tabular}
\end{table}

%======================================================================
\section{工程集成说明}
%======================================================================

\subsection{文件结构}

\begin{lstlisting}
20260626-6.5.6-调制方式/
  slb_phy_mod.h
  slb_phy_mod.c
  SLB物理层6.5.6调制方式技术文档.tex
  SLB_6.5.6_调制方式_学习资料.tex
\end{lstlisting}

\subsection{与标准其他章节的衔接}

\begin{itemize}[nosep]
  \item \textbf{6.5.7}：数据调制前须完成加扰，输入为 $\tilde{g}_k$；
  \item \textbf{6.5.3}：参考信号 QPSK 独立实现，不得调用本模块；
  \item \textbf{6.5.4}：调制符号映射到 RE 后叠功控；
  \item \textbf{6.2.6}：编码速率与调制阶数共同决定 MCS。
\end{itemize}

\subsection{编译}

\begin{lstlisting}[language=bash]
cc -std=c99 -Wall -Wextra -c slb_phy_mod.c -o slb_phy_mod.o
cc your_phy_module.c slb_phy_mod.o -lm -o your_phy_app
\end{lstlisting}

\vfill
\begin{center}
\small\textcolor{gray}{精确参数与字段定义以 T/XS 10001-2025 正式标准为准；本文档仅供 SLB 物理层实现与联调参考。}
\end{center}

\end{document}
